% Generated draft: requires mathematical and source review.
\documentclass[11pt]{article}
\usepackage[margin=25mm]{geometry}
\usepackage{fontspec}
\setmainfont{Latin Modern Roman}
\usepackage{amsmath,amssymb,mathtools,unicode-math}
\setmathfont{Latin Modern Math}
\usepackage{enumitem,needspace,xurl,hyperref}
\setlength{\parindent}{0pt}
\setlength{\parskip}{5pt}
\setlength{\emergencystretch}{3em}
\newcommand{\N}{\mathbb{N}}
\newcommand{\Z}{\mathbb{Z}}
\newcommand{\Q}{\mathbb{Q}}
\newcommand{\R}{\mathbb{R}}
\newcommand{\C}{\mathbb{C}}
\newcommand{\F}{\mathbb{F}}
\newcommand{\A}{\mathbb{A}}
\newcommand{\PP}{\mathbb{P}}
\newcommand{\E}{\mathbb{E}}
\newcommand{\eps}{\varepsilon}
\newcommand{\dd}{\,\mathrm{d}}
\begin{document}
\section*{UnsolvedMath ID: 20003344}
\textbf{Source problem code:} catalogue entry
\subsection{Definitions and mathematical statement}
This is the first part of a literature-status sentence following B. Weiss's Conjecture 48, not a standalone conjecture. It says that [200] formulated the conjecture and checked special cases, then stops before saying which ones. The following companion record names one-dimensional $D$ [38] and the full maximal real split torus case $D=A$ [188].

The antecedent setup, in companion permanent ID $20003342$, has a semisimple real algebraic group $G$ with arithmetic structure, a noncocompact arithmetic lattice $\Gamma$, and a closed subgroup $D$ of a maximal real split torus $A$. For $x=g\Gamma$, divergence means properness of $D\to G/\Gamma$, $d\mapsto dx$: every sequence leaving compact subsets of $D$ has image leaving compact subsets of $G/\Gamma$. The source calls divergence ``obvious'' when finitely many open subsemigroups covering $D$ have representation-vector contraction certificates. For the usual rational certificate, the representations are defined over $\Q$ and the witness vectors must be nonzero rational vectors, with the arithmetic lattice structure ensuring discreteness; the terse antecedent does not explicitly state all those witness conditions. Allowing a zero witness would make the criterion vacuous.

Conjecture 48 compares $\dim D$ with $\operatorname{rank}_{\Q}G$: no divergent orbits above the rational rank, only obvious ones at equality, and nonobvious ones below it. This is a historical formulation. Companion research material warns that the last clause needs almost-$\Q$-simple or appropriate factorwise scope and fails literally for arbitrary semisimple products, so a blanket modern resolution of all three clauses is not asserted.

The present aid supplies a precise active-weight criterion for a single sector. Let $W$ be a finite-dimensional real exponent space, $C\subset W$ a closed, pointed, full-dimensional convex cone, and $S=\operatorname{int}C$. Let a real split torus representation have weight decomposition $V=\bigoplus_\chi V_\chi$, and let $0\ne w=\sum_\chi w_\chi$. A weight is active if $w_\chi\ne0$. Write
\[
 C^*=\{\ell\in W^*: \ell(H)\ge0\text{ for every }H\in C\}.
\]
The reported equivalence is
\[
 \rho(\exp H_j)w\longrightarrow0
 \text{ for every }H_j\in S\text{ with }\|H_j\|\to\infty
 \quad\Longleftrightarrow\quad
 -\chi\in\operatorname{int}(C^*)\text{ for every active }\chi.
\]
Equivalently there are $K,\varepsilon>0$ such that
\[
 \|\rho(\exp H)w\|\le Ke^{-\varepsilon\|H\|}\quad(H\in S).
\]
Thus strict contraction on every unbounded sector sequence requires a uniform margin including directions approaching its boundary. For an irredundantly faceted polyhedral cone $C=\{H:\ell_j(H)\ge0,\ 1\le j\le m\}$, the finite certificate is
\[
 -\chi=\sum_{j=1}^m c_{\chi j}\ell_j,\qquad c_{\chi j}>0
 \quad\text{for every active }\chi.
\]
Interiors here are taken in the displayed exponent space; the full-dimensional convention avoids silently using an ambient interior for a lower-dimensional sector.

For the modular lattice model $a_t=\operatorname{diag}(e^t,e^{-t})$, a frame $g$ whose two inverse coordinate lines are rational has nonzero $v_+,v_-\in\Q^2$ with $gv_+\in\R e_2$, $gv_-\in\R e_1$. These are explicit standard-representation shrinking witnesses on the forward and backward sectors, respectively; this includes $g\in\mathrm{SL}_2(\Q)$. The criterion concerns rational certificates for such models, not existence of a certificate for every divergent orbit. All auxiliary results are supplied machine-generated aid requiring independent verification.
\subsection{Short English statement}
This fragment reports special cases of Weiss's divergent-orbit conjecture. Its supplied mathematical content is a sector test: every active weight must be strictly negative with a uniform margin on the closed cone, giving a finite positive-normal certificate for polyhedral sectors.
\subsection{Sources}
\begin{itemize}
\item B. Weiss's Conjecture 48 and source reference [200], American Institute of Mathematics, \emph{Emerging applications of measure rigidity}, section 8, \url{https://aimath.org/WWN/measrigid/measrigid.pdf}. The following source sentence cites [38] for a one-dimensional case and [188] for $D=A$; full bibliography entries are absent.
\item Companion records in \texttt{aim-other-notes.json}: ID $20003342$, \texttt{AIM-OTHER-0149}, \texttt{notes[148]}, supplies the appended divergence and obviousness setup; ID $20003343$, \texttt{AIM-OTHER-0150}, \texttt{notes[149]}, states the three clauses and its aid flags the product-group correction; ID $20003345$, \texttt{AIM-OTHER-0152}, \texttt{notes[151]}, supplies the special-case continuation.
\item Local UnsolvedMath registry, permanent ID $20003344$, source code \texttt{catalogue entry}, canonical \texttt{aim-other-notes.json}, \texttt{notes[150]}, dated 14 August 2026. Research attempt 1: context only, ``not a problem'' at run, proved rigor, low novelty confidence; machine-generated and not peer reviewed.
\end{itemize}
\end{document}
